\section{Legal and Institutional Framework}
\label{sec:institutional}

\paragraph{Procurement statutes.}
Brazilian public procurement during our 2009--2019 sample is
governed by Lei~8.666/93 (the General Procurement Statute) and Lei
10.520/2002 (Pregão). Two modalities dominate the BEC data we
analyze. \emph{Convite} (Article~22 of Lei~8.666) is a sealed-bid
invitation procedure: the buyer invites at least three suppliers,
all proposals are submitted in writing, and the contract is awarded
to the qualifying lowest bid. The statute requires three valid
proposals for the procedure to clear; if fewer appear, the buyer
must either re-issue the call or document the absence of bidders.
\emph{Pregão} (Lei~10.520) is a reverse electronic auction: bidders
submit progressively lower offers in real time until no firm
improves the standing best offer. Pregão has no minimum-bidder
requirement.

The structural difference between the two modalities is central
to our analysis. The convite minimum-bidder rule was introduced
to ensure procedural competition; in practice, when a buyer's
short-list yields fewer than three motivated bidders, additional
proposals must be obtained. The competition-economics literature
\citep{marshall2012economics, asker2010leniency,
porter1993detection} identifies this configuration---a rule
mandating $\underline{n}$ bidders, but no rule that those bidders
genuinely compete---as the canonical setting for
\emph{regulatory cover bidding}: designated firms submit
deliberately losing bids to satisfy the procedural requirement
while preserving the rents of the chosen winner.

The 2021 Brazilian reform (Lei~14.133/2021) replaces Lei~8.666 for
new contracts, restricts the convite modality, and tightens
habilitation requirements. The reform postdates our sample but
shapes the policy relevance of the FL screen: even where statutory
incentives for cover bidding are reduced, residual minimum-bidder
provisions in subnational systems and analogous provisions in OECD
procurement regimes preserve the incentive structure the screen
exploits.

\paragraph{Enforcement architecture.}
Three institutions police BEC procurement. The
\textbf{Tribunal de Contas do Estado de S\~ao Paulo} (TCE-SP)
audits purchasing units administratively, with binding power to
require correction of irregularities and report findings to
prosecutorial bodies. The \textbf{Controladoria Geral do Estado}
(CGE) oversees executive-branch compliance with procurement law.
The \textbf{Conselho Administrativo de Defesa Econ\^omica} (CADE)
is the federal competition authority; it investigates and
adjudicates cartel and bid-rigging cases under Lei~12.529/2011,
imposing administrative fines (typically 0.1\%--20\% of firm
revenue) and may refer cases to the Public Prosecutor's Office for
criminal prosecution under Lei~8.137/90 (Article~3).

CADE's procurement-cartel portfolio during our window comprises
12 adjudicated cases covering 65 firm-defendants, with final
decisions spanning 2015--2025. Eight of the twelve cases were
decided after our 2019 sample cut, so portions of our external
validation are prospective: 2009--2019 FL classification
anticipates firms whose cartel involvement enforcement confirmed
years later. The interval between operational cartel activity and
final CADE adjudication---the cases in our portfolio span 4--16
years from process opening to ruling---is important context for
the validation design and for the policy claim that screening at
the buyer level can shorten this interval substantially.

\paragraph{The screening problem.}
Bid-level forensic tools---variance, kurtosis, and spread analyses
across submitted bids---detect coordinated bidding when the
required microdata are preserved
\citep{imhof2019detecting, huber2019machine, wallimann2023machine}.
Many electronic procurement systems, including BEC during much of
our sample, archive only winning prices and participant lists at
the contract-record layer; bid-by-bid records may exist in
operational logs but are not always exported into the analytical
warehouses that oversight bodies query. A screen that requires
only participation records lowers the data threshold for proactive
detection to the level virtually every system maintains. That
data-availability argument is the screen's primary contribution
to the law-and-economics of procurement enforcement: not a more
powerful test, but a test deployable where the more powerful
tests are not.
